% ---------------------------------------------------------------------
%  Chapitre 1 — Introduction : but + planification de l'expérience
%  (couvre les points « but du stage » et « planification » des consignes)
% ---------------------------------------------------------------------
\chapter{Introduction}
\label{chap:intro}

\section{But du mémoire}
\label{sec:but}

% ANGLE DU MÉMOIRE (décidé) : ce n'est PAS une étude de stratégies de
% trading, mais la construction d'un OUTIL DE VALIDATION statistique.
% Les stratégies ne sont que des jeux de données de test.

Lors de mon stage en entreprise j'ai eu la consigne de créer un agent IA capable de faire du trading.
J'ai alors décidé de donner à mon agent des capacités d'analyse statistique.
Après avoir créé le nécessaire pour qu'un agent puisse passer des ordres et
avoir accès à différentes données du marché, j'ai estimé qu'il fallait qu'il soit capable d'évaluer
les stratégies de trading qu'il aura établies.

Je vais présenter dans ce rapport l'outil d'évaluation de stratégies que
j'ai donné à mon agent.

L'objectif de cet outil n'est pas de découvrir une stratégie de trading
rentable. Il est de construire un enchaînement de méthodes capable de
déterminer, pour une stratégie quelconque qu'on lui soumet, si l'avantage
qu'elle affiche est réel ou s'il relève du hasard d'échantillonnage
ou du biais de marché.

Autrement dit, la contribution revendiquée est méthodologique~:
la chaîne de tests, sa justification et la vérification de ses conditions
d'application. 

Les dix stratégies d'analyse technique utilisées dans ce
rapport ne constituent pas l'objet d'étude mais un banc d'essai~:
elles fournissent des jeux de données réels permettant de
montrer que l'outil discrimine correctement — qu'il sait rejeter
ce qui ne tient pas et retenir ce qui résiste à l'analyse. Les résultats
obtenus sont donc présentés à titre d'illustration du bon
fonctionnement du protocole, et non comme des recommandations
d'investissement.

Pour s'assurer que l'outil ne rejette pas tout par excès de sévérité,
on lui soumettra aussi un témoin dont l'avantage est réel par construction.


\section{Les trois menaces qui pèsent sur le jugement}
\label{sec:menaces}

% [À RÉDIGER — le cœur de l'introduction]
% Idées à développer avec TES mots (annonce les trois obstacles que le
% protocole devra franchir ; chacun sera traité dans le rapport) :
%  1. LA DÉRIVE DU MARCHÉ : sur la période, la quasi-totalité des titres
%     montent -> un rendement positif ne prouve aucun talent. (Déjà
%     amorcé en §\ref{sec:but} : réponse = l'edge apparié au hasard.)
%  2. LA MULTIPLICITÉ : dix stratégies testées d'un coup -> multiplier
%     les occasions d'un faux positif (~40 % de risque global à 5 %).
%     Réponse = contrôle du risque global (Bonferroni).
%  3. LA DÉPENDANCE : les trades ne sont pas des témoins indépendants
%     (même titre, mêmes journées de bourse, chevauchements). C'est la
%     menace la plus grave et la moins visible -> elle sera MESURÉE
%     avant tout test (chap. 2) puis neutralisée pas à pas (chap. 3).
%  - Finir sur l'idée : une validation qui ignore une seule de ces
%    menaces produit des verdicts trop optimistes ; l'outil les traite
%    toutes les trois, et c'est sa raison d'être.

La première difficulté tient à la dérive du marché.
Le premier réflexe pour juger une stratégie est d'examiner si elle
est rentable : un gain positif la ferait « marcher ». Ce raisonnement
est pourtant trompeur, car encore faut-il que ce gain dépasse ce
qu'une entrée au hasard aurait rapporté.

Ici les stratégies sont appliquées aux titres du S\&P~500,
l'indice des 500 grandes capitalisations américaines,
dont la tendance est haussière sur la période étudiée.
Cette dérive haussière est un problème car une stratégie peut effectuer
des gains sans pourtant n'avoir capté aucun signal.
C'est pourquoi on calculera non pas les gains mais plutôt l'avantage
qu'apporte chaque stratégie.

La deuxième provient du nombre de tests. En effet,
plus on teste de stratégies, plus le risque d'obtenir au moins un faux
positif augmente. Ce problème sera traité au
chapitre~\ref{sec:b1-etape1} à l'aide de la correction de Bonferroni.

La troisième, sans doute la plus délicate, est la dépendance
des données. De nombreuses méthodes statistiques
imposent en condition l'indépendance des variables, or ici les trades sont
dépendants sous différents aspects, comme la période, le titre et même le secteur~;
cette difficulté sera mesurée au chapitre~\ref{chap:donnees}, puis neutralisée
au chapitre~\ref{chap:juger}.

On est forcé de prendre en compte ces trois difficultés, sous peine
d'un verdict trop optimiste : c'est leur traitement conjoint qui fait de ce
protocole un véritable outil de validation.

\section{Démarche du mémoire}
\label{sec:planification}

% [À RÉDIGER — en PROSE, 1 à 2 paragraphes, pas une liste à puces]
% Idée directrice : UNE seule enquête, du diagnostic au verdict. Enchaîner
% les 4 chapitres avec des transitions, dans TES mots. Squelette :
%
%  - Phrase d'ouverture (garde l'esprit) : ce mémoire suit une seule
%    enquête, du diagnostic au verdict.
%  - Ch. 2 (\ref{chap:donnees}) : on présente les données (les 40 362
%    trades, l'edge apparié) PUIS on mesure la dépendance qui menace les
%    tests -> on part du terrain et de ses pièges.
%  - Ch. 3 (\ref{chap:juger}) : le CŒUR. On juge les stratégies, du test
%    naïf au test corrigé (l'escalade de prudence), puis on lit le verdict
%    à la lumière du marché lui-même.
%  - Ch. 4 (\ref{chap:contexte}) : on change de question -> comparer les
%    stratégies entre elles et caractériser le contexte d'entrée (ANOVA,
%    khi-deux, analyses factorielles). Statut descriptif.
%  - Ch. 5 (\ref{chap:protocole}) : on récapitule l'outil en livrable,
%    ses limites et les étapes suivantes.
%  - Optionnel (1 phrase) : dire que le fil va du plus fragile (la donnée,
%    la dépendance) au plus solide (le verdict, puis l'outil réutilisable).
%
% Astuce transitions : « d'abord… », « le cœur du travail… », « on change
% alors de question… », « enfin… ». Utilise les \ref{} pour les numéros.

L'étude suit le fil d'une enquête unique, du diagnostic au verdict.

% NOTE (ne pas rédiger) : l'ancienne architecture « par blocs » (stratégies
% / signaux d'information / relations inter-actifs / régression finale) a été
% resserrée : le mémoire se concentre sur le jugement des stratégies ; les
% signaux d'information et la régression sont renvoyés en perspectives
% (chap.~\ref{chap:protocole}).
